% ====================================================================
\section{The Proposed Trust Model}
\label{sec:proposedModel}
% ====================================================================
This section presents the three extensions to the IoDT
model~\cite{soula2024iodt}, whose architecture is shown in
Fig.~\ref{fig:architecture}.  Twin replicas run Extension~I, a quorum
over their signed votes (Extension~II) decides on the real-time path,
and verdicts and votes are committed to a signed Merkle audit log off
that path (Extension~III).

\Figure[!t](topskip=0pt,botskip=0pt,midskip=0pt)[width=\textwidth]{figs/fig_architecture.png}
{Architecture of the proposed IoDT trust pipeline with its three extensions.\label{fig:architecture}}

\subsection{Extension I: Continuous Trust Score}
\label{subsec:ext1-math}
The four factors are mapped to $[0,1]$ and combined with non-negative
weights that sum to one:
\begin{equation}
s_{t}(i) = w_{E}E_{t}(i) + w_{P}\frac{P_{t}(i)}{100}
         + w_{T}T_{t}(i) + w_{W}W_{t}(i).
\label{eq:weighted_score}
\end{equation}
All experiments use $w_{E}=w_{P}=0.30$ and $w_{T}=w_{W}=0.20$, which
put slightly more weight on the physical factors than on the
notification-driven ones.  A single bad round should not condemn a
node that has otherwise behaved well, so the score is smoothed by an
exponential moving average,
\begin{equation}
\tau_{t}(i) = \alpha\,s_{t}(i) + (1-\alpha)\,\tau_{t-1}(i),
\qquad \tau_{0}(i)=1,
\label{eq:ema}
\end{equation}
and the node is flagged when $\tau_t(i)<\theta$.  The defaults are
$\alpha=0.5$ and $\theta=0.70$, selected on a separate tuning seed
(Section~\ref{sec:setup}).

Smoothing has a blind spot: a weighted score cannot react to an
attacker that always fails the \emph{same} low-weight factor.  An
attacker that fails only $W$ keeps $s_t\in[0.776,0.80]$, and, absent
congestion, its EMA never falls below $0.70$.  Extension~I therefore adds a
\emph{per-factor persistence check}: the node is also flagged when any
single factor has failed in at least $m$ of the last $w$ rounds.
Congestion scatters failures across factors and rounds, whereas a
liar fails the same factor again and again.  The defaults are $m=4$
and $w=8$.  A second, longer window ($m_2$ of $w_2$ rounds) can be
added against rule-aware attackers (Section~\ref{sec:res-stealth}); it
is off by default.  Algorithm~\ref{alg:ema} gives the complete
per-round update.

\begin{algorithm}[!t]
\caption{Extension~I trust update for node $i$ at round $t$.}
\label{alg:ema}
\noindent\textbf{Input:} factor scores $(E_t,P_t,T_t,W_t)$; weights;
prior trust $\tau_{t-1}(i)$; failure counters of the last $w$ (and
$w_2$) rounds; $\alpha$, $\theta$, $m$, $w$ (and $m_2$, $w_2$).

\noindent\textbf{Procedure:}
\begin{enumerate}\setlength{\itemsep}{0pt}\setlength{\parskip}{0pt}
  \item $s_t \gets w_E E_t + w_P (P_t/100) + w_T T_t + w_W W_t$.
  \item $\tau_t(i) \gets \alpha\,s_t + (1-\alpha)\,\tau_{t-1}(i)$.
  \item Update, for each factor, the number of failures in rounds
        $t-w+1,\dots,t$ (a factor fails when it is $0$, or when
        $P_t<\Theta_{\mathrm{PDR}}$).
  \item If $\tau_t(i)<\theta$, or some factor failed in $\ge m$ of the
        last $w$ rounds (or in $\ge m_2$ of the last $w_2$), return
        \textsc{flag}; else return \textsc{trust}.
\end{enumerate}
\end{algorithm}

\subsubsection{False-positive analysis}
The analysis rests on three assumptions, and the experiments test what
happens when the first two are relaxed.
\begin{itemize}
\item[(A1)] \emph{Independence.}  For an honest node, factor
corruptions are independent across factors and rounds.
\item[(A2)] \emph{Symmetric corruption.}  Each factor is corrupted
with probability $q$; a corrupted binary factor becomes $0$ and a
corrupted PDR falls below $\Theta_{\mathrm{PDR}}$.
\item[(A3)] \emph{Healthy-link PDR.}  An uncorrupted PDR lies in
$[\underline{P},100]$; the simulator uses $\underline{P}=92\%$.
\end{itemize}
Under (A1)--(A2) the binary AND (\ref{eq:binary_and}) flags an honest
node whenever any factor is corrupted, so its per-round FPR is
\begin{equation}
\mathrm{FPR}_{\mathrm{AND}}(q) = 1-(1-q)^{4} = 4q + O(q^{2}),
\label{eq:fpr_and}
\end{equation}
and by the same counting any rule that flags only after two failures
--- the AND required to fail in two consecutive rounds, or two of the
four factors required to fail in the same round --- errs at rate
$\Theta(q^{2})$.  In the evaluation the 2-of-4 rule is combined with
the persistence check below, which keeps its order at $\Theta(q^{2})$
and lets it see a liar that fails only one factor.  These rules serve as delay-matched baselines in
Section~\ref{sec:resultsAndDiscussion}.  The EMA needs three failures.

\begin{proposition}[FPR of the EMA score]\label{prop:fpr_ema}
A node is flagged by the EMA at round $t$ if and only if
\begin{equation}
s_t < s^{\star}=\frac{\theta-(1-\alpha)\,\tau_{t-1}}{\alpha}.
\label{eq:sstar}
\end{equation}
With the default weights, $\alpha=0.5$ and $\theta=0.70$, a flag
requires at least three factor failures within the last $L=6$ rounds,
so under (A1)--(A3)
\begin{equation}
\mathrm{FPR}_{\mathrm{EMA}}(q) \le \binom{4L}{3}q^{3},\qquad
\mathrm{FPR}_{\mathrm{EMA}}(q)=\Theta(q^{3}).
\end{equation}
\end{proposition}
\begin{proof}
See Appendix~\ref{app:proofs}.
\end{proof}

The cubic order depends on the parameters and on (A1).  At
$\alpha\geq0.6$ (with $\theta=0.70$) or $\theta\geq0.75$ (with
$\alpha=0.5$), two failures in one round can already cross the
threshold and the order falls to $\Theta(q^{2})$
(Appendix~\ref{app:proofs}).  When one congestion event fails several
factors of a node at once, (A1) fails across factors, and the order is
again $\Theta(q^{2})$ (Section~\ref{sec:res-robust}).  The price of
smoothing is one round of delay: without congestion, a node that lies
in every round scores $s_t<0.60$, so from $\tau_0=1$ its EMA is below
$0.70$ after two rounds, whereas the AND flags it in the first.

The persistence check adds false positives at a controlled rate and
catches lies that the EMA cannot see.

\begin{proposition}[Persistence check]\label{prop:persist}
Under (A1)--(A2) the probability that the persistence check flags an
honest node in a given round is at most
$4\binom{w}{m}q^{m}$ ($7.2\times10^{-4}$ at $m=4$, $w=8$, $q=0.04$).
A node whose lies fail a given factor independently in each round with
probability $p$ is flagged within any window of $w$ rounds with
probability at least $\Pr[\mathrm{Bin}(w,p)\geq m]$; for $p=1$ it is
flagged by round $m$.  A node that keeps at most $m-1$ lies in every
window of $w$ rounds is never flagged by the check: its \emph{lie
budget} is $(m-1)/w$, or $\min\{(m-1)/w,(m_2-1)/w_2\}$ with a second
window.
\end{proposition}
\begin{proof}
See Appendix~\ref{app:proofs}.
\end{proof}

% ====================================================================
\subsection{Extension II: Multi-Twin Voting}
\label{sec:ext2}
% ====================================================================
A single twin is a single point of trust: if it is corrupted, the
node it watches is never flagged.  Extension~II deploys $n$ replicas
of each twin and decides by a quorum of their signed votes; the
reported configuration is $n=5$, $k=3$.  This is one-shot voting over
signed verdicts in the sense of Byzantine quorum
systems~\cite{malkhi1998byzantine}, not a full Byzantine-fault-tolerant
agreement protocol such as PBFT~\cite{castro2002practical}: there is no
multi-phase commit, view change or total order on operations.  The
extension relies on the following assumptions.
\begin{itemize}
\item[(B1)] \emph{Independent hosting.}  The replicas run on
administratively independent hosts (different fog sites or cloud
availability zones), so compromising one host does not compromise the
others.
\item[(B2)] \emph{Authenticated, bound votes.}  In every round each
replica signs one vector of verdicts for all the nodes it watches,
bound to its identity, the round number and the current key epoch
(ECDSA over secp256k1 in our implementation).  Keys are rotated by
starting a new epoch, which invalidates signatures from retired keys.
\item[(B3)] \emph{Bounded delay.}  Votes from honest replicas arrive
within a known bound $\Delta$; a vote that arrives later is treated as
missing.
\item[(B4)] \emph{Identical evidence.}  The fog node forwards the same
evidence to every replica, so honest replicas reach identical
verdicts.
\item[(B5)] \emph{Adversary budget.}  At most $c\le k-1$ replicas are
corrupted.
\end{itemize}
Each replica runs Algorithm~\ref{alg:ema}, signs its vote vector and
sends it to the trust engine, which verifies the signatures and
applies the \emph{quorum rule} of Algorithm~\ref{alg:vote}:
\textsc{trust} if at least $k$ valid votes say \textsc{trust},
\textsc{flag} if at least $n-k+1$ say \textsc{flag}, and otherwise the
previous verdict is kept and an availability alarm raised.
The engine forwards the verdict to the fog node, which enforces it
and, being trusted (Table~\ref{tab:threats}), appends the verdict it
enforced to the off-chain store of Extension~III.  With all
$n=2k-1$ votes present exactly one quorum is reached, so the rule
coincides with ``flag if fewer than $k$ \textsc{trust} votes''; unlike
that rule, missing or withheld votes never cause a new isolation,
although a node that is already isolated stays isolated while no
quorum forms.  The count is computed by the trust engine, not by a smart
contract, so no ledger operation lies on the decision path.  The
verdict and the signed votes are appended to the off-chain store of
Extension~III.

\begin{algorithm}[!t]
\caption{Per-round vote for node $i$.}
\label{alg:vote}
\noindent\textbf{Input:} evidence $\mathbf{x}_t(i)$; replicas
$\mathrm{DT}_1,\dots,\mathrm{DT}_n$; quorum $k$; previous verdict.

\noindent\textbf{Procedure:}
\begin{enumerate}\setlength{\itemsep}{0pt}\setlength{\parskip}{0pt}
  \item Each $\mathrm{DT}_r$ runs Algorithm~\ref{alg:ema} on
        $\mathbf{x}_t(i)$ and signs its vote vector ($\sigma_r$),
        bound to $(r,t,\text{epoch})$.
  \item The trust engine waits at most $\Delta$, verifies each received
        signature and its binding, and counts the \textsc{trust} votes
        $V_T$ and the \textsc{flag} votes $V_F$.
  \item If $V_T\geq k$ return \textsc{trust}; if $V_F\geq n-k+1$ return
        \textsc{flag} to the fog node, which isolates node $i$;
        otherwise keep the previous verdict and raise an alarm.
  \item Append the verdict and the signed votes to the off-chain store
        (input to Algorithm~\ref{alg:anchor}).
\end{enumerate}
\end{algorithm}

Under (B3)--(B5) the quorum behaves like one honest twin.  Corrupted
replicas that always vote \textsc{trust} for an attacker
(\emph{shielding}) are outvoted whenever the $n-c\geq n-k+1$ honest
replicas flag it, and corrupted replicas that vote \textsc{flag} for an
honest node (\emph{framing}) are outvoted whenever the honest replicas
trust it.  Detection and false positives are therefore those of an
honest single twin for every $c\le k-1$, whereas a single corrupted
twin never flags its node and $c\geq k$ corrupted replicas decide every
verdict.

The trust engine itself is not replicated.  It is trusted to count
promptly, but not trusted for integrity: every enforced verdict is
logged with the signed votes behind it, and replicas endorse a header
of the audit log only if each verdict in it follows from its votes by
the quorum rule (Section~\ref{sec:ext3}).  An engine that issues an
unsupported verdict is therefore exposed when the next header is
committed, within $K$ rounds plus the time to append it, although it
is not prevented in real time.  Voting defends
against corrupted twins; it does not help against attackers that lie
rarely, whose detection is governed by
Eq.~(\ref{eq:sstar}) and Proposition~\ref{prop:persist}.
If the replicas receive independent evidence instead of (B4), framing
replicas can tip honest nodes whose honest replicas disagree
(Section~\ref{sec:res-ext2}).

% ====================================================================
\subsection{Extension III: Signed Merkle Audit Log}
\label{sec:ext3}
% ====================================================================
Most blockchain-based trust schemes consult or update the chain while
a decision is being made (Section~\ref{sec:res-comparison}).  Here the
verdict is produced by Extensions~I and~II alone, and an audit log
records it afterwards.  The log is a hash-chained sequence of signed
headers over Merkle-committed off-chain data, in the tradition of
tamper-evident logs~\cite{schneier1999secure,crosby2009efficient,laurie2013rfc6962}.
Its headers can be kept in any append-only store; hosting them on a
permissioned ledger (for example Hyperledger Fabric with a
crash-tolerant ordering service among fog peers) adds replicated
storage and availability against an operator that deletes the log, but
the tamper evidence of Theorem~\ref{thm:tamper} does not depend on the
ledger's consensus.

\emph{Content.}  Every round, the per-node event (factor scores,
$\tau_t$, verdict; $|E|=29$~bytes) and the replicas' signed vote
vectors are stored off-chain.  Every $K$ rounds the trust engine
proposes a header that carries the height, the first and last round,
the numbers of events and votes, the key epoch, the previous header
hash, and the Merkle roots of the window's events, of its signed votes
and of the per-node state records ($|R|=13$~bytes per node);
$|H|=152$~bytes.  The trees follow RFC~6962~\cite{laurie2013rfc6962}:
leaf and interior hashes are domain-separated, the last node is not
duplicated, and the leaf count is committed in the header.  Each
replica \emph{endorses} the header --- signs it --- only after checking
the off-chain data against it, including that every verdict follows
from its votes; a header needs $k$ endorsements.  A partial final
window is committed as well.  Each header therefore occupies
$|H|+k\,|\sigma|$ bytes in the header store, with $|\sigma|\le73$~bytes per
endorsement (identifier and DER signature), independent of the number
of nodes.  Raw evidence never
leaves the off-chain store, which also limits what the log reveals
about each node.  A verdict becomes auditable within $K$ rounds plus
the time to append a header; this delay never affects when the verdict
itself is issued.  We further assume:
\begin{itemize}
\item[(C1)] header hashes are periodically published on an
independently controlled channel (linked
time-stamping~\cite{haber1991time}, a KSI hash
tree~\cite{buldas2013ksi}, an Arbitrum-style state
commitment~\cite{kalodner2018arbitrum}, or a public-chain anchor);
\item[(C2)] the signature scheme is existentially unforgeable and
the hash function is collision resistant.
\end{itemize}

\begin{algorithm}[!t]
\caption{Audit-log commit at round $t$.}
\label{alg:anchor}
\noindent\textbf{Input:} pending events $\mathcal{E}$ and votes
$\mathcal{S}$; state records; period $K$; previous header hash
$H_{\mathrm{prev}}$.

\noindent\textbf{Procedure:}
\begin{enumerate}\setlength{\itemsep}{0pt}\setlength{\parskip}{0pt}
  \item Append the round's events and signed votes (off-chain).
  \item If $t \bmod K \neq 0$ and the run is not ending, return.
  \item The trust engine forms the header with
        $\mathrm{MTH}(\mathcal{E})$, $\mathrm{MTH}(\mathcal{S})$,
        $\mathrm{MTH}(\text{state})$, counts, epoch and
        $H_{\mathrm{prev}}$.
  \item Each replica checks the window (lengths, counts, roots, vote
        signatures and bindings, verdicts against votes, state against
        events) and, if it passes, signs the header.
  \item With $k$ endorsements, append the header; periodically
        publish its hash externally (C1).
\end{enumerate}
\end{algorithm}

\begin{definition}[Verifier]
Given the headers and the off-chain data, the verifier $V$ checks the
hash links, at least $k$ valid endorsements from distinct replicas of
the header's key epoch, record lengths and counts, every Merkle root,
every vote signature and its binding to replica, round and epoch, that
each verdict follows from its round's votes by the quorum rule, that
each state record equals the node's latest event, and, if a checkpoint
$(j,h^{\star}_j)$ is available, that header $j$ hashes to $h^{\star}_j$.
\end{definition}

\begin{theorem}[Tamper evidence]\label{thm:tamper}
Under (C2), let the adversary modify, insert, delete or reorder any
event, vote or state record of a committed window $B_i$ and recompute
any hashes it likes.
\begin{enumerate}
\item If it holds fewer than $k$ replica signing keys, $V$ rejects
except with negligible probability.
\item If it holds $k$ or more keys, it can make $V$ accept if and only
if no checkpoint $(j,h^{\star}_j)$ with $j\ge i$ was published before
the modification (up to a hash collision).
\end{enumerate}
\end{theorem}
\begin{proof}
See Appendix~\ref{app:proofs}.
\end{proof}

Endorsements, not the checkpoint, are therefore the first line of
defence; the checkpoint is what defeats an adversary that steals or
colludes with $k$ replicas.

\subsection{End-to-End Round}
\label{sec:round}
Fig.~\ref{fig:sequence} shows how the three extensions interact in one
round.  Steps 1--7 form the real-time decision path: the fog node
forwards the evidence to the replicas, each replica runs
Algorithm~\ref{alg:ema} and signs its vote vector, and the trust engine
applies Algorithm~\ref{alg:vote} and returns the verdict to the fog
node, which isolates a flagged node.  Steps 8--11 run afterwards and
never delay a verdict: events and votes are appended to the off-chain
store, and every $K$ rounds Algorithm~\ref{alg:anchor} forms a header,
collects $k$ replica endorsements and appends it; header hashes are
periodically published on the checkpoint channel (C1).

\Figure[!t](topskip=0pt,botskip=0pt,midskip=0pt)[width=0.92\textwidth]{figs/fig_sequence.png}
{Message sequence of one round: the decision path (steps 1--7) and the off-path audit-log commit (steps 8--11).\label{fig:sequence}}

